\documentclass[10pt]{article}

% The documentation loads the library itself: every example is compiled live.
\usepackage[seed=2026]{fancytests}
\usepackage[T1]{fontenc}
\usepackage{lmodern}
\renewcommand*\familydefault{\sfdefault}
\geometry{a4paper, top=22mm, bottom=22mm, left=22mm, right=22mm}

\usepackage[most]{tcolorbox}
\usepackage{hyperref}
\hypersetup{colorlinks, linkcolor=FriendlyDenim, urlcolor=FriendlyDenim}

% ---------------------------------------------------------------------------
% Documentation helpers
% ---------------------------------------------------------------------------
\lstdefinestyle{ftcode}{
  basicstyle=\ttfamily\footnotesize,
  keywordstyle=\color{FriendlyDenim},
  commentstyle=\color{black!50},
  columns=fullflexible, keepspaces=true, breaklines=true,
  literate={à}{{\`a}}1 {è}{{\`e}}1 {é}{{\'e}}1 {ì}{{\`i}}1 {ò}{{\`o}}1 {ù}{{\`u}}1
           {’}{{'}}1,
}
\lstset{style=ftcode}

% code above, result below
\newtcblisting{example}{
  enhanced, colback=white, colframe=black!20, colbacklower=white,
  listing above text, sharp corners, boxrule=.5pt,
  segmentation style={black!20, solid},
  listing options={style=ftcode},
  before upper={\ResetAllScores\testsetup{subject=math}},
  title={\small Example}, coltitle=black!70, colbacktitle=black!4,
  fonttitle=\bfseries}
% code only
\newtcblisting{code}{
  enhanced, colback=black!3, colframe=black!12, listing only, sharp corners,
  boxrule=.4pt, listing options={style=ftcode}}

\newcommand*{\cs}[1]{\texttt{\textbackslash#1}}
\newcommand*{\env}[1]{\texttt{#1}}
\newcommand*{\key}[1]{\texttt{\color{FriendlyDenim}#1}}
\newcommand*{\meta}[1]{\ensuremath{\langle}\textit{#1}\ensuremath{\rangle}}
\newcommand*{\opt}[1]{\texttt{#1}}
\newcommand*{\marg}[1]{\texttt{\{}\meta{#1}\texttt{\}}}
\newcommand*{\oarg}[1]{\texttt{[}\meta{#1}\texttt{]}}

% key tables
\newenvironment{keys}
  {\par\smallskip\noindent\small
   \renewcommand{\arraystretch}{1.25}%
   \tabularx{\linewidth}{@{}>{\raggedright\arraybackslash}p{.24\linewidth}
                          >{\raggedright\arraybackslash}p{.17\linewidth}>{\raggedright\arraybackslash}X@{}}
   \toprule \textbf{Key} & \textbf{Default} & \textbf{Meaning}\\ \midrule}
  {\bottomrule\endtabularx\par\smallskip}

\newenvironment{cmds}
  {\par\smallskip\noindent\small
   \renewcommand{\arraystretch}{1.25}%
   \tabularx{\linewidth}{@{}>{\raggedright\arraybackslash}p{.36\linewidth}>{\raggedright\arraybackslash}X@{}}
   \toprule \textbf{Command} & \textbf{Meaning}\\ \midrule}
  {\bottomrule\endtabularx\par\smallskip}

\newcommand*{\note}[1]{%
  \par\smallskip\noindent
  \colorbox{FriendlyYellow!18}{\parbox{\dimexpr\linewidth-2\fboxsep}{\small\textbf{Note.} #1}}%
  \par\smallskip}

% the test header switches the page style off: keep page numbers in the manual
\NewCommandCopy\ftOrigTestheader\testheader
\RenewDocumentCommand\testheader{O{}}{\ftOrigTestheader[#1]\pagestyle{plain}}

\setlength{\parskip}{.45em}
\setlength{\emergencystretch}{3em}
\linespread{1.1}

\title{\Huge\bfseries fancytests\\[.3em]
  \Large\mdseries Maths and physics tests in \LaTeX}
\author{Matteo Bramardi\\[.2em]
  \normalsize co-authored with Claude (Anthropic)}
\date{Version 1.0 --- October 2026}

\begin{document}
\pagestyle{plain}
\maketitle
\thispagestyle{empty}

\begin{abstract}
\noindent
\textsf{fancytests} is a \LaTeX{} package for writing school tests in
mathematics and physics. It replaces the test-related part of
\textsf{fancylessons}, the author's previous package, with a smaller and
faster one. Every command takes
named options (\texttt{key=value}); there are no more round, angle or bar
delimiters.

It provides: a test header (full or compact) with the boxes required for
students with DSA; exercises and questions with automatic score totals;
multiple-choice questions with shuffled options and an exported answer key;
reproducible randomisation of variants, numbers, questions and exercises;
reduced versions of a test from the same source; teacher copies with
solutions; standard styles for function graphs and physics diagrams.

An optional compatibility layer compiles the tests written with
\textsf{fancylessons} without changes.
\end{abstract}

\tableofcontents
\newpage

%% ===========================================================================
\section{Getting started}
%% ===========================================================================

\subsection{About this package}

\textsf{fancylessons} was the author's earlier \LaTeX{} package, used both
for lesson notes (theorem and definition boxes, colours, page styles) and for
tests. Over time it grew to more than 4000 lines and about 70 loaded packages,
which made every compilation slow. \textsf{fancytests} keeps only what is
needed for tests, redesigned with a uniform \texttt{key=value} syntax and
new features (answer key, shuffled questions, reduced versions, teacher
copies).

The references to \textsf{fancylessons} in this manual (speed comparisons,
the \opt{legacy} option, the table of old and new commands in
section~\ref{sec:legacy}) are there mainly for the author, to compare the old
and the new version and to reuse past tests. Other readers can safely ignore
them.

\textsf{fancytests} and this manual were written by Matteo Bramardi together
with Claude, an AI model by Anthropic, which wrote and tested most of the
code under the author's direction.

\subsection{Installation}

The package is a single file, \texttt{fancytests.sty}. Put it in the same
folder as the document and load it in the preamble.

On Overleaf, the simplest way to keep it up to date is to link it from GitHub:
\emph{Upload} $\to$ \emph{From external URL}, with the address
\begin{center}\small
\texttt{https://raw.githubusercontent.com/Bradwave/fancytests/main/fancytests.sty}
\end{center}
The \emph{Refresh} button of the linked file downloads the latest version.
Replace \texttt{main} with a release tag (e.g.\ \texttt{v1.0}) to freeze a
project on that version, which is advisable for tests already given.

\note{On Overleaf, use \TeX{} Live 2023 or later (Menu $\to$ Settings $\to$
\TeX{} Live version). Older versions lack some \LaTeX3 functions used by the
package.}

\subsection{A minimal test}

\begin{code}
\documentclass[10pt]{article}
\usepackage{fancytests}
\usepackage[T1]{fontenc}
\usepackage{lmodern}
\renewcommand*\familydefault{\sfdefault}
\geometry{a4paper, top=15mm, bottom=15mm, left=17.5mm, right=17.5mm}
\linespread{1.25}

\testsetup{class=4M, subject=math, dsa=rid}

\begin{document}
\testheader[date=20/03/2026, title={Logaritmi}]

\exercise Risolvi le seguenti equazioni.
\begin{parts}[cols=2]
  \itemx[pts=1] $3^x = 9$
  \itemx[pts=1] $\log_2 x = 3$
\end{parts}

\exercise Rispondi.
\begin{parts}
  \itemx[pts=1] Quale funzione è sempre crescente?
    \begin{choices}
      \correct $e^x$
      \choice  $e^{-x}$
      \choice  $\ln(-x)$
      \choice[fixed] Nessuna delle precedenti
    \end{choices}
\end{parts}

\vfill
\finalscore
\end{document}
\end{code}

The package already loads \textsf{babel} (Italian), \textsf{geometry},
\textsf{amsmath}, \textsf{amssymb}, \textsf{mathtools}, \textsf{tabularx},
\textsf{booktabs}, \textsf{enumitem}, \textsf{xcolor}, \textsf{bm},
\textsf{cancel}, \textsf{tikz}, and, unless disabled, \textsf{pgfplots} and
\textsf{siunitx}. Do not load them again with different options.

\subsection{Package options}

Options are given when loading the package:
\verb|\usepackage[seed=42, answers]{fancytests}|.

\begin{keys}
\key{subject} & \opt{math} & Default subject: \opt{math} or \opt{physics}.
  Sets the colour (blue or green) and the word in the title.\\
\key{plots} & \opt{true} & Load \textsf{pgfplots}. With \opt{false} the
  \env{ftplot} environment is not available, but loading is faster.\\
\key{units} & \opt{true} & Load \textsf{siunitx} (for \cs{qty}, \cs{num}, \ldots),
  set up with the decimal comma.\\
\key{seed} & \opt{1} & Seed of the random generator (section~\ref{sec:random}).
  \opt{0} = a different seed at every compilation.\\
\key{answers} & off & Teacher copy: correct options highlighted, \env{solution}
  boxes printed.\\
\key{nofigures} & off & Fast draft: every TikZ picture and graph is replaced by
  an empty frame.\\
\key{cache} & off & Store figures and recompile them only when their code
  changes (package \textsf{memoize}). Needs shell escape, which is enabled on
  Overleaf.\\
\key{legacy} & off & Load the old \textsf{fancylessons} commands
  (section~\ref{sec:legacy}).\\
\end{keys}

\subsection{Compilation speed}

On a typical test the package loads in about $0.8$\,s, against $1.9$\,s for
\textsf{fancylessons}. Most of the remaining time is spent drawing graphs.
To speed things up:

\begin{itemize}[nosep]
\item while writing, use \opt{nofigures}: pictures become placeholders;
\item use \opt{cache}: each picture is drawn once and reused until its code
  changes. The first two compilations are slower, the following ones much
  faster;
\item if a test has no graphs, use \opt{plots=false}.
\end{itemize}

\note{With \opt{cache}, a picture whose content depends on randomisation
through a macro defined elsewhere (e.g.\ a shuffled graph stored in a macro)
cannot be recognised as changed. After changing such pictures or the seed,
use \emph{Recompile from scratch} on Overleaf, which empties the cache.}

%% ===========================================================================
\section{Global settings: \cs{testsetup}}
%% ===========================================================================

\cs{testsetup}\marg{keys} sets options for everything that follows. In the
preamble it sets the defaults of the document; the same keys can be given to
\cs{testheader} for a single test.

\begin{keys}
\key{subject} & \opt{math} & \opt{math} / \opt{physics} (also \opt{matematica},
  \opt{fisica}); any other text is printed as given.\\
\key{class} & --- & Class, e.g.\ \opt{4M}.\\
\key{date} & empty & Date; when empty, a blank \_\_/\_\_/\_\_\_\_ is printed.\\
\key{title} & empty & Subtitle of the test (topic).\\
\key{type} & \opt{Verifica} & Type of test: \opt{Verifica}, \opt{Test},
  \opt{Interrogazione}, \ldots\\
\key{style} & \opt{full} & Header layout: \opt{full} or \opt{compact}.\\
\key{titlebreak} & \opt{auto} & Compact header only: put the title on its own
  line when it does not fit next to the subject (\opt{auto}), always
  (\opt{always}) or never, letting it wrap (\opt{never}).\\
\key{dsa} & \opt{none} & DSA boxes: \opt{none}, \opt{dsa}, \opt{rid}
  (section~\ref{sec:header}).\\
\key{timing} & \opt{true} & Print start, end and hand-in time fields.\\
\key{recovery} & empty & Date of the original test, for a resit
  (``recupero del \ldots'').\\
\key{gradebox} & \opt{VOTO} & Label of the grade box.\\
\key{showversion} & off & Print the copy number next to the class (\#1, \#2,
  \ldots).\\
\key{reduced} & off & Reduced version: questions marked \key{noreduced} are
  left out.\\
\key{answers} & off & Same as the package option, but can be switched on and
  off inside the document.\\
\key{labels} & \opt{box} & Labels of multiple-choice options: \opt{box}
  ($\square$) or \opt{letters} (A, B, C, \ldots).\\
\key{shuffle} & \opt{true} & Shuffle multiple-choice options.\\
\key{shuffleparts} & \opt{false} & Shuffle the questions of every \env{parts}.\\
\key{seed} & \opt{1} & Restart the random generator from this seed.\\
\key{decimal} & \opt{.} & Decimal mark in the score boxes.\\
\key{opacity} & \opt{12} & Intensity (\%) of the coloured score boxes.\\
\key{grade-min}, \key{grade-max} & \opt{2}, \opt{10} & Lowest and highest grade
  in the grade formula.\\
\key{final} & \opt{compact} & Default style of \cs{finalscore}.\\
\key{autoclose} & \opt{true} & Close automatically an exercise left open
  (section~\ref{sec:exercises}).\\
\end{keys}

The header keys also accept Italian names: \key{materia}, \key{classe},
\key{data}, \key{titolo}, \key{tipo}, \key{orari}, \key{recupero}, \key{voto}.

%% ===========================================================================
\section{The header}\label{sec:header}
%% ===========================================================================

\cs{testheader}\oarg{keys} starts a new test: it prints the header, resets
exercise numbers and scores and increases the copy number. It accepts all the
keys of \cs{testsetup}; the values given here hold for this test only.

\subsection{Full header}

\begin{example}
\testheader[class=4M, date=20/03/2026,
  title={Logaritmo, equazioni esponenziali e logaritmiche}]
\end{example}

\subsection{Compact header}

\key{style=compact} uses about half the height. The title and the class share
the first row, names and times share the second, the DSA boxes and the resit
date go on an optional third row. A thin rule closes the header. When the
title does not fit on the first row it moves, as a whole, to a row of its own
(key \key{titlebreak}).

\begin{example}
\testheader[style=compact, class=4M, subject=physics,
  title={Onde e acustica}, dsa=dsa, recovery=12/05/2026]
\end{example}

\subsection{DSA boxes, times, resits}

\begin{itemize}[nosep]
\item \key{dsa=dsa}: \emph{Schemi e calcolatrice}, \emph{Riduzione es./tempo
  aggiuntivo}, \emph{Compensazione};
\item \key{dsa=rid}: \emph{Uso della calcolatrice}, \emph{Uso di schemi/mappe},
  \emph{Riduzione es./tempo aggiuntivo};
\item \key{dsa=none}: no boxes.
\end{itemize}
\key{timing=false} removes the time fields (e.g.\ for oral tests);
\key{recovery=}\meta{date} prints ``recupero del \meta{date}''.

\begin{example}
\testheader[class=4G, type=Interrogazione, subject=physics,
  timing=false, date=02/02/2026]
\end{example}

%% ===========================================================================
\section{Exercises}\label{sec:exercises}
%% ===========================================================================

\cs{exercise}\oarg{keys} prints the heading ``Esercizio $n$'' followed by the
text of the exercise.

\begin{keys}
\key{diff} & --- & Difficulty from 1 to 3, shown as coloured squares.\\
\key{title} & --- & A different heading, e.g.\ \opt{title=Test intermedi}.
  \cs{exercise}\verb|[Test intermedi]| also works.\\
\key{bonus} & off & Bonus exercise: heading ``Esercizio Bonus*'', points go to
  the bonus, not to the total. \cs{bonusexercise} is a shortcut.\\
\key{fixed}, \key{follow} & --- & Used only when exercises are shuffled
  (section~\ref{sec:shuffle}).\\
\end{keys}

\subsection{Closing an exercise}

\cs{exercisescore}\oarg{keys} closes the exercise and prints its score box on
the right. Without options the score is the sum of the questions. Options:
\key{pts} imposes the score (\cs{exercisescore}\verb|[5]| is the same as
\verb|[pts=5]|), \key{label} changes the label (default E1, E2, \ldots),
\key{short} makes the box narrower. For bonus exercises use
\cs{bonusexercisescore}; \cs{bonusexercisenote} prints the explanatory note.

If an exercise is not closed, the next \cs{exercise} or \cs{finalscore} closes
it (key \key{autoclose}). Close it explicitly before a \cs{newpage}, otherwise
the score box ends up at the top of the next page.

\begin{example}
\exercise[diff=2] Un'onda si propaga con velocità $v = \qty{340}{m/s}$.
\begin{parts}
  \itemx[pts=1] Calcola la lunghezza d'onda per $f = \qty{170}{Hz}$.
  \itemx[pts=0.5] Come varia $\lambda$ se $f$ raddoppia?
\end{parts}
\exercisescore

\exercise Spiega la differenza tra scala lineare e logaritmica.
\exercisescore[2]

\bonusexercise Determina il tempo di dimezzamento dello Iodio-131.
\bonusexercisescore[0.5]
\bonusexercisenote
\end{example}

%% ===========================================================================
\section{Questions: the \env{parts} environment}
%% ===========================================================================

Questions are written inside \env{parts}, one \cs{itemx} each. Letters
continue across several \env{parts} of the same exercise and restart at every
\cs{exercise}. The score of each question appears at the right of its first
line.

\subsection{Options of \env{parts}}

\begin{keys}
\key{stretch} & \opt{1.5} & Spacing between questions, as a multiple of the
  line height (\opt{1} = no extra space).\\
\key{cols} & \opt{1} & Number of columns; questions fill the rows from left to
  right.\\
\key{start} & --- & Letter to start from (\opt{start=3} starts at ``c'').\\
\key{ragged} & off & Ragged-right text instead of justified.\\
\key{shuffle} & off & Shuffle the questions (section~\ref{sec:shuffle}).\\
\end{keys}

\subsection{Options of \cs{itemx}}

\begin{keys}
\key{pts} & --- & Points of the question; added to the exercise.\\
\key{diff} & --- & Difficulty 1--3, squares at the left of the letter.\\
\key{mark} & off & Highlight the letter.\\
\key{upper} & off & Capital letter.\\
\key{noreduced} & off & Leave the question out of the reduced version.\\
\key{fixed}, \key{follow} & --- & For shuffled questions.\\
\end{keys}

\begin{example}
\exercise Risolvi.
\begin{parts}[cols=2, stretch=1.3]
  \itemx[pts=1] $3^x - 2^x = 0$
  \itemx[pts=1] $\log_{10}(2x-1) + \log_{10}(x) = 1$
  \itemx[pts=1, diff=2] $e^{2x} - e^{x+1} = 0$
  \itemx[pts=1, mark] $5 \cdot 2^x - 12 = 2^x$
\end{parts}
\end{example}

\subsection{Graphs and questions side by side}

\env{parts} works inside a \env{minipage}; the scores stay at the right edge
of the minipage.

\begin{example}
\exercise Osserva il grafico.

\begin{minipage}[c]{.36\linewidth}
  \centering
  \begin{ftplot}[xmin=-1.9, xmax=8.1, ymin=-3.9, ymax=3.1,
                 width=5.5cm, height=5cm]
    \fplot[domain=0.01:8]{ln(x)/ln(3)}
  \end{ftplot}
\end{minipage}\hfill
\begin{minipage}[c]{.6\linewidth}
  \begin{parts}[ragged]
    \itemx[pts=1] Quale funzione è rappresentata?
    \itemx[pts=1] Determina dominio e immagine.
  \end{parts}
\end{minipage}
\end{example}

\subsection{Reduced versions}

Mark with \key{noreduced} the questions to leave out of the reduced (DSA)
version, and compile with \cs{testsetup}\verb|{reduced}|. The omitted
questions disappear, the following letters move up and their points are not
counted. Full and reduced versions come from the same file.

\begin{example}
\testsetup{reduced}
\exercise Rispondi.
\begin{parts}
  \itemx[pts=1] Primo quesito.
  \itemx[pts=2, noreduced] Quesito omesso nella versione ridotta.
  \itemx[pts=1] Terzo quesito, stampato come b.
\end{parts}
\exercisescore
\testsetup{reduced=false}
\end{example}

%% ===========================================================================
\section{Scores and grade}
%% ===========================================================================

\subsection{Score commands}

\begin{cmds}
\cs{score}\oarg{short}\marg{n} & Add $n$ points to the current exercise and
  print the box. Inside \env{parts} use \key{pts} instead.\\
\cs{invscore}\marg{n} & Add points without printing anything.\\
\cs{scorebox}\oarg{short}\marg{n} & Print a box without adding points.\\
\cs{emptyscorebox} & An empty box.\\
\cs{exercisescore}\oarg{keys} & Close the exercise (section~\ref{sec:exercises}).\\
\cs{currenttotal} & Total of the closed exercises so far.\\
\cs{ResetAllScores} & Reset everything (done by \cs{testheader}).\\
\end{cmds}

Scores are computed with exact decimal arithmetic and printed with up to two
decimals: $1 \to 1.0$, $0.75 \to 0.75$.

\subsection{Final score: \cs{finalscore}}

\cs{finalscore}\oarg{keys} closes the last exercise, prints the summary and
resets the scores, ready for the next copy.

\begin{keys}
\key{style} & \opt{compact} & \opt{table}: table E1, E2, \ldots, EB, Tot.\ and
  formula; \opt{compact}: formula and total on one line; \opt{formula}: formula
  only; \opt{minimal}: total only; \opt{none}.\\
\key{max} & sum & Maximum score shown, if different from the computed total.\\
\key{bonus} & \opt{auto} & Bonus column in the table: \opt{auto} (only if there
  are bonus points), \opt{true}, \opt{false}.\\
\key{label} & \opt{Tot:} & Label of the total.\\
\end{keys}

The formula is
$\text{voto} = \frac{\text{punti ottenuti}}{\text{punti massimi}}
\cdot (\text{max} - \text{min}) + \text{min}$ with \key{grade-min} and
\key{grade-max} of \cs{testsetup} (default $2$ and $10$).

\begin{example}
\exercise Primo. \begin{parts}\itemx[pts=2] a \itemx[pts=1.5] b\end{parts}
\exercise Secondo. \begin{parts}\itemx[pts=3] c\end{parts}
\bonusexercise Bonus. \bonusexercisescore[0.5]
\finalscore[style=table]
\end{example}

\begin{example}
\exercise Unico. \exercisescore[6]
\finalscore
\end{example}

%% ===========================================================================
\section{Multiple-choice questions}
%% ===========================================================================

\subsection{The \env{choices} environment}

Inside \env{choices}, \cs{correct} starts the correct option and \cs{choice}
a wrong one. Options are shuffled; the number of columns is chosen
automatically: all on one row if they fit, otherwise two columns, otherwise
one. Long options wrap with a hanging indent.

\begin{example}
\exercise Rispondi.
\begin{parts}
  \itemx[pts=1] Quale funzione ha dominio $x \ge -3$?
    \begin{choices}
      \correct $\sqrt{x+3}$
      \choice  $\log(x+3)$
      \choice  $e^{x+3}$
      \choice  $\frac{1}{x+3}$
    \end{choices}
  \itemx[pts=1] Se $f(x)$ è una decrescita esponenziale, $\ln f(x)$ è
    \begin{choices}
      \correct una retta con pendenza negativa.
      \choice  una retta con pendenza positiva.
      \choice  una curva logaritmica decrescente.
      \choice[fixed] nessuna delle precedenti.
    \end{choices}
\end{parts}
\end{example}

\subsection{Options}

\begin{keys}
\key{cols} & \opt{auto} & \opt{auto}, or a fixed number of columns.\\
\key{shuffle} & \opt{true} & Shuffle the options (global default in
  \cs{testsetup}).\\
\key{labels} & \opt{box} & \opt{box} or \opt{letters}.\\
\key{indent} & \opt{0pt} & Left indent (\cs{itmindent} outside \env{parts}).\\
\key{before} & \opt{2pt} & Space above the options.\\
\key{rowsep} & \opt{0pt} & Extra space between rows of options.\\
\key{after} & \opt{.25em} & Space below the options.\\
\end{keys}

Options of a single choice:
\cs{choice}\verb|[fixed]| keeps the option in its position (typically the last
one, ``Nessuna delle precedenti''); \cs{correct}\verb|[if={v=2}]| makes the
option correct only when the variant \opt{v} chosen by \cs{pick} is the second
(section~\ref{sec:random}).

Options can contain anything, including TikZ pictures.

\subsection{Answer key and teacher copy}

Every compilation writes the file \meta{document name}\texttt{.key.csv}, with
one line per multiple-choice question:

\begin{code}
versione;esercizio;quesito;risposta
1;2;b;C
1;2;c;A
2;2;b;B
\end{code}

\texttt{versione} is the copy number (one per \cs{testheader}),
\texttt{esercizio} the exercise number (\texttt{B} for the bonus),
\texttt{quesito} the letter of the question, \texttt{risposta} the letter of
the correct option in the printed order. On Overleaf the file is in
\emph{Logs and output files}.

With the option \opt{answers} (or \cs{testsetup}\verb|{answers}|) the correct
options are highlighted. Since the seed is the same, the teacher copy matches
the student copies exactly.

\begin{example}
\testsetup{answers, labels=letters}
\exercise Rispondi.
\begin{parts}
  \itemx[pts=1] Capitale d'Italia?
    \begin{choices}
      \correct Roma \choice Milano \choice Torino \choice[fixed] Nessuna
    \end{choices}
\end{parts}
\begin{solution}
Roma è la capitale dal 1871.
\end{solution}
\testsetup{answers=false, labels=box}
\end{example}

\subsection{Solutions}

The \env{solution}\oarg{title} environment is printed only in the teacher copy
(\opt{answers}); otherwise it is ignored. The default title is
\emph{Soluzione}.

%% ===========================================================================
\section{Randomisation}\label{sec:random}
%% ===========================================================================

\subsection{The seed}

All random choices (shuffled options, questions and exercises, \cs{pick},
\cs{randint}, the old \textsf{fancylessons} commands) come from one
generator. Its seed is fixed, so every compilation gives the same copies and
the teacher copy and the answer key always match.

\begin{cmds}
\verb|\usepackage[seed=42]{fancytests}| & Seed of the whole document
  (default \opt{1}).\\
\cs{testsetup}\verb|{seed=42}| & Same; inside the document, it restarts the
  generator from that point.\\
\cs{setseed}\verb|{42}| & Restart the generator from 42 here (e.g.\ before each
  class).\\
\verb|seed=0| & A different seed at every compilation.\\
\end{cmds}

\note{Do not use \opt{seed=0} for a real test: the answer key and the teacher
copy would only match the student copies of the same compilation. To get new
variants, change the number.}

\subsection{Variants and numbers}

\begin{cmds}
\cs{pick}\oarg{name}\verb|{a \or b \or c}| & One of the variants at random. With
  a name, the choice is remembered for \cs{correct}\verb|[if={name=2}]|.\\
\cs{randint}\oarg{name}\marg{min}\marg{max} & A random integer between min
  and max, remembered with the name.\\
\cs{useval}\marg{name} & Print again the number drawn with that name.\\
\end{cmds}

\begin{example}
\exercise
Per quale funzione \pick[d]{$\Dom f = (1, +\infty)$ \or $\Dom f = (-1, +\infty)$}
si ha $\Ima f = \spR$?
\begin{choices}
  \correct[if={d=1}] $f(x) = -\ln(x - 1)$
  \correct[if={d=2}] $f(x) = -\ln(x + 1)$
  \choice $f(x) = -e^{x - 1}$
  \choice $f(x) = e^{x + 1}$
\end{choices}

Calcola $\randint[a]{2}{9} \cdot \randint[b]{2}{9}$.
Il risultato è $\useval{a} \cdot \useval{b} = \inteval{\useval{a}*\useval{b}}$.
\end{example}

\subsection{Shuffling blocks}

The \env{shuffle}\oarg{name} environment shuffles a list of blocks, which are
then printed with \cs{shuffled}\marg{name}\marg{position}. Useful for graphs or
pictures to be placed in a random order.

\begin{example}
\begin{shuffle}[g]
  \block{\tikz\draw[thick] (0,0) -- (1,1);}
  \block{\tikz\draw[thick] (0,1) -- (1,0);}
  \block{\tikz\draw[thick] (0,.5) -- (1,.5);}
\end{shuffle}
(a) \shuffled{g}{1} \qquad (b) \shuffled{g}{2} \qquad (c) \shuffled{g}{3}
\end{example}

\subsection{Shuffling questions and exercises}\label{sec:shuffle}

\env{parts}\verb|[shuffle]| shuffles its questions (\cs{testsetup}%
\verb|{shuffleparts}| does it for every \env{parts}). The environment
\env{exercises} shuffles the exercises it contains. Letters and exercise
numbers are assigned after shuffling; scores and the answer key follow.

\begin{itemize}[nosep]
\item \key{fixed} keeps a question or exercise in its position (counted among
  the blocks: see below);
\item \key{follow} ties it to the previous one, so they move together.
\end{itemize}

\begin{example}
\exercise Rispondi.
\begin{parts}[shuffle]
  \itemx[pts=1] Uno.
  \itemx[pts=1] Due.
  \itemx[pts=1] Tre: osserva il grafico.
  \itemx[pts=1, follow] Quattro: con riferimento al quesito precedente\dots
  \itemx[pts=1, fixed] Cinque (sempre nella quinta posizione tra i blocchi).
\end{parts}
\end{example}

\begin{code}
\begin{exercises}
  \exercise ...            % each exercise runs until the next \exercise
  \exercise[fixed] ...     % stays in its place
  \exercise[follow] ...    % moves together with the previous one
\end{exercises}
\end{code}

Limitations:
\begin{itemize}[nosep]
\item a question tied with \key{follow} forms a block with the previous one;
  \key{fixed} keeps the position among blocks, so letters after a block of two
  move down by one;
\item shuffled questions must not refer to each other by letter;
\item inside \env{exercises}, \cs{exercise} must not be inside a
  \env{minipage}; keep the bonus exercise outside the environment;
\item \cs{verb} cannot be used inside \env{parts}.
\end{itemize}

\subsection{Several copies}

\cs{Repeat}\marg{n}\marg{code} repeats a test $n$ times. Each \cs{testheader}
starts a new copy with new random choices; \key{showversion} prints the copy
number, which is also the \texttt{versione} column of the answer key.

\begin{code}
\newcommand{\verifica}{\input{docs/3-logaritmo/base}\newpage}
\Repeat{6}{\verifica}
\end{code}

%% ===========================================================================
\section{Graphs}
%% ===========================================================================

The \env{ftplot}\oarg{axis options} environment draws a function graph with
the standard style \opt{ft}: centred axes, major and minor grid, thick black
curves. Successive curves change style (solid, dashed, grey, dotted) so they
can be told apart in black-and-white photocopies.

\cs{fplot}\oarg{options}\marg{expression}\oarg{legend} adds a curve; the
options are those of \cs{addplot} (\opt{domain}, \opt{samples}, \opt{dashed},
\ldots). Any \textsf{pgfplots} option can be passed to \env{ftplot}.

\begin{example}
\begin{center}
\begin{ftplot}[xmin=-4.9, xmax=5.1, ymin=-4.9, ymax=5.1, width=7cm, height=7cm]
  \fplot[domain=-5:5]{e^x - 1}[A]
  \fplot[domain=1.001:5]{-ln(x - 1)}[B]
  \fplot[domain=-0.999:5]{ln(x + 1)}[C]
  \fplot[domain=-5:-0.001]{ln(-x)}[D]
\end{ftplot}
\end{center}
\end{example}

The style can also be used in an ordinary \env{axis}:
\verb|\begin{axis}[ft, xmode=log]|. The default number of samples is $80$;
increase it with \opt{samples=150} for curves with sharp bends.

%% ===========================================================================
\section{Physics diagrams}
%% ===========================================================================

TikZ styles for physics diagrams:

\begin{keys}
\key{force} & --- & Thick red arrow.\\
\key{velocity} & --- & Thick blue arrow.\\
\key{vector} & --- & Generic black arrow.\\
\key{axis arrow} & --- & Thin arrow for axes.\\
\key{dimension} & --- & Two-headed arrow for lengths.\\
\key{body} & --- & Outlined body with light fill (for nodes).\\
\key{ground} & --- & Grey fill for floors, slopes, walls.\\
\key{spring} & --- & Spring (coil).\\
\key{continue}, \key{startcontinue}, \key{endcontinue} & \opt{.75em} &
  Dotted continuation at both ends, at the start, at the end of a curve.\\
\end{keys}

\begin{example}
\begin{center}
\begin{tikzpicture}
  \fill[ground] (0,0) -- (4,0) -- (4,2) -- cycle;
  \draw[thick] (0,0) -- (4,0) -- (4,2) -- cycle;
  \node[body, minimum size=6mm, rotate=26.57] (b) at (2.2,1.45) {};
  \draw[force] (b.center) -- ++(0,-1.1) node[right] {$m\vb{g}$};
  \draw[velocity] (b.north) -- ++(-0.9,-0.45) node[above] {$\vb{v}$};
  \draw[dimension] (0,-.3) -- node[below] {$L$} (4,-.3);
  \draw[spring] (5,0) -- (5,2);
  \draw[thick, continue, domain=-1:1, samples=40] plot ({\x+7}, {\x*\x});
\end{tikzpicture}
\end{center}
\end{example}

%% ===========================================================================
\section{Mathematics, units, text}
%% ===========================================================================

\begin{cmds}
\cs{spN}, \cs{spZ}, \cs{spQ}, \cs{spR}\oarg{exp}, \cs{spC}\oarg{exp} &
  $\spN, \spZ, \spQ, \spR, \spR[2], \spC$\\
\cs{Dom}, \cs{Ima}, \cs{Sign} & $\Dom f$, $\Ima f$, $\Sign x$\\
\cs{sen}, \cs{tg}, \cs{arcsen}, \cs{arctg} & $\sen x$, $\tg x$, $\arcsen x$,
  $\arctg x$\\
\cs{abs}\marg{x}, \cs{norm}\marg{x} & $\abs{x}$, $\norm{x}$\\
\cs{eqdef}, \cs{eqq} & $\eqdef$, $\eqq$\\
\cs{lesq}, \cs{gesq} & $\lesq$, $\gesq$\\
\cs{vb}\marg{F}, \cs{vect}\marg{AB} & $\vb{F}$, $\vect{AB}$\\
\cs{compacteq}, \cs{quitecompacteq}, \cs{notsocompacteq} & Reduce the space
  around the next displayed equations (3, 6, 9\,pt).\\
\cs{unm}\marg{u}, \cs{sunm}\oarg{exp}\marg{u}\marg{v} & Light units:
  $3\unm{kg}$, $9.8\sunm[2]{m}{s}$. For full units use \textsf{siunitx}:
  \cs{qty}\verb|{340}{m/s}| gives \qty{340}{m/s}.\\
\midrule
\cs{tbf}, \cs{tit}, \cs{tbfit}, \cs{trm} & Bold, italic, bold italic, roman.\\
\cs{red}, \cs{bred}, \cs{blue}, \cs{bblue} & \red{red}, \bred{bold red},
  \blue{blue}, \bblue{bold blue}.\\
\cs{hlc}\oarg{colour}\marg{text}, \cs{highlight}\marg{text} & \hlc{grey
  background}, \highlight{subject colour}.\\
\cs{dotline}\marg{length} & A dotted line to fill in: \dotline{6em}\\
\end{cmds}

Colours: \textcolor{MathsBlue}{\texttt{MathsBlue}},
\textcolor{PhysicsGreen}{\texttt{PhysicsGreen}}, \texttt{SubjectColor} (the
colour of the current subject) and the colour-blind friendly palette
\textcolor{FriendlyRed}{\texttt{FriendlyRed}},
\textcolor{FriendlyOrange}{\texttt{FriendlyOrange}},
\textcolor{FriendlyYellow}{\texttt{FriendlyYellow}},
\textcolor{FriendlyGreen}{\texttt{FriendlyGreen}},
\textcolor{FriendlyForest}{\texttt{FriendlyForest}},
\textcolor{FriendlyBlue}{\texttt{FriendlyBlue}},
\textcolor{FriendlyDenim}{\texttt{FriendlyDenim}},
\textcolor{FriendlyMidnight}{\texttt{FriendlyMidnight}}.

%% ===========================================================================
\section{Compatibility with fancylessons}\label{sec:legacy}
%% ===========================================================================

With \verb|\usepackage[legacy]{fancytests}| the tests written with
\textsf{fancylessons} compile without changes. The old commands are read only
when the option is given, so they do not slow down new tests. In legacy mode
exercises are not closed automatically, as in \textsf{fancylessons}.

\begin{small}
\renewcommand{\arraystretch}{1.25}
\noindent
\begin{tabularx}{\linewidth}{@{}>{\raggedright\arraybackslash}p{.47\linewidth}X@{}}
\toprule
\textbf{fancylessons} & \textbf{fancytests}\\ \midrule
\cs{mathtestheader}\verb|[rec]{4M}{date}{title}(dsa)| &
  \cs{testheader}\verb|[class=4M, date=..., title=..., dsa=dsa]|\\
\cs{phystestheader}, \cs{physinttestheader}, \ldots &
  \key{subject=physics}, \key{type=Interrogazione}, \key{timing=false}\\
\verb|\begin{scoredlist}(1.5)<L>| \ldots\ \verb|& \score{1}\\| &
  \verb|\begin{parts}[ragged]| \ldots\ \cs{itemx}\verb|[pts=1]|\\
\env{splitlist}, \env{threesplitlist} & \env{parts}\verb|[cols=2]|,
  \verb|[cols=3]|\\
\cs{randomfouroptions}\verb|{a}{b}{c}{d}(.5)<L>|\verb+|0pt|+ &
  \verb|\begin{choices} \correct a \choice b ...|\\
\cs{randomfouroptionsline}, \ldots\verb|list|, \ldots\verb|plusone|\ldots &
  \env{choices}\verb|[cols=4]|, \verb|[cols=1]|, \cs{choice}\verb|[fixed]|\\
\cs{randOp}\verb|{a}{b}|, \cs{randOpThree}, \cs{randSign} &
  \cs{pick}\verb|{a \or b}|, \cs{pick}\verb|{+ \or -}|\\
\cs{randIntR}\verb|{1}{9}|, \cs{randInt} & \cs{randint}\verb|{1}{9}|\\
\cs{pgfmathdeclarerandomlist} + \cs{prunelist} & \env{shuffle} +
  \cs{shuffled}\\
\cs{minimalfinalscore} + \cs{finalscorenotable} &
  \cs{finalscore}\verb|[style=compact]|\\
\cs{finalscorenotable} & \cs{finalscore}\verb|[style=formula]|\\
\cs{overrideAllColors}\marg{c} & \cs{colorlet}\verb|{SubjectColor}{c}|\\
\bottomrule
\end{tabularx}
\end{small}

Also available in legacy mode: \env{tflist}, \env{splittflist},
\env{tflistinv}, \cs{truefalse}, \cs{truefalseinv}, \cs{twooptions},
\cs{fouroptions}, \cs{fouroptionslist}, \cs{fiveoptionslist},
\cs{randomthreeoptionsline}, \cs{randomfiveoptions}(\texttt{line},
\texttt{list}), \cs{randomfourplusoneoptionslist}, \cs{mathscore},
\cs{physscore}, \cs{intermediatetest}, \cs{lengthof}, \cs{nthof},
\env{notex*}, \cs{mathsol}, \cs{physsol} and the unit shortcuts
\cs{unmeter}, \cs{unmpers}, \ldots

\note{The old option macros do not know which option is correct, so their
questions do not appear in the answer key. To use the key, rewrite them with
\env{choices} and \cs{correct}.}

%% ===========================================================================
\section{A complete test}
%% ===========================================================================

\begin{code}
\documentclass[10pt]{article}
\usepackage[seed=2026]{fancytests}
\usepackage[T1]{fontenc}
\usepackage{lmodern}
\renewcommand*\familydefault{\sfdefault}
\geometry{a4paper, top=15mm, bottom=15mm, left=17.5mm, right=17.5mm}
\linespread{1.25}

\testsetup{class=4M, dsa=none, style=compact, showversion}

\newcommand{\verifica}{%
  \testheader[date=20/03/2026, title={Logaritmi}]

  \exercise Risolvi le seguenti equazioni.
  \begin{parts}[cols=2, shuffle]
    \itemx[pts=1] $3^x - 2^x = 0$
    \itemx[pts=1] $\log_{10}(2x-1) + \log_{10}(x) = 1$
    \itemx[pts=1, noreduced] $2^x + 2 = 0$
    \itemx[pts=1, noreduced] $5 \cdot 2^x - 12 = 2^x$
  \end{parts}

  \exercise Rispondi.
  \begin{parts}
    \itemx[pts=1] Quale funzione è sempre crescente?
      \begin{choices}
        \correct $e^x$ \choice $e^{-x}$ \choice $-\ln x$
        \choice[fixed] Nessuna
      \end{choices}
  \end{parts}
  \begin{solution} $e^x$ è crescente perché la base è maggiore di 1. \end{solution}

  \vfill
  \finalscore
  \newpage}

\begin{document}
\Repeat{4}{\verifica}          % 4 different copies
% \testsetup{reduced}          % reduced (DSA) version
% \usepackage[answers]{...}    % teacher copy: same seed, same copies
\end{document}
\end{code}

\end{document}
